% Generated draft: requires mathematical and source review.
\documentclass[11pt]{article}
\usepackage[margin=25mm]{geometry}
\usepackage{fontspec}
\setmainfont{Latin Modern Roman}
\usepackage{amsmath,amssymb,mathtools,unicode-math}
\setmathfont{Latin Modern Math}
\usepackage{enumitem,needspace,xurl,hyperref}
\setlength{\parindent}{0pt}
\setlength{\parskip}{5pt}
\setlength{\emergencystretch}{3em}
\newcommand{\N}{\mathbb{N}}
\newcommand{\Z}{\mathbb{Z}}
\newcommand{\Q}{\mathbb{Q}}
\newcommand{\R}{\mathbb{R}}
\newcommand{\C}{\mathbb{C}}
\newcommand{\F}{\mathbb{F}}
\newcommand{\A}{\mathbb{A}}
\newcommand{\PP}{\mathbb{P}}
\newcommand{\E}{\mathbb{E}}
\newcommand{\eps}{\varepsilon}
\newcommand{\dd}{\,\mathrm{d}}
\begin{document}
\section*{UnsolvedMath ID: 30003723}
\textbf{Source problem code:} catalogue entry
\subsection{Definitions and mathematical statement}
Use the nearest-neighbour graph on \(\Z^2\). Adopt the height convention in which a stable configuration is \(\eta\in\{0,1,2,3\}^{\Z^2}\), and a toppling at a vertex removes four grains there and sends one grain to each of its four neighbours. Let \(\nu\) denote the stationary measure referred to in the supplied record; its precise construction is not supplied.

Fix an addition site \(x\in\Z^2\). To describe the avalanche without assuming that infinite-volume stabilization terminates, let \(V_n=\{-n,\ldots,n\}^2\), with \(n\) sufficiently large that \(x\in V_n\). Restrict \(\eta\) to \(V_n\), add one grain at \(x\), and stabilize with all grains leaving \(V_n\) absorbed by a sink. Let \(u_{n,\eta}^{x}(y)\) be the number of topplings at \(y\), extended by zero outside \(V_n\), and define the extended toppling count
\[
 u_{\eta}^{x}(y):=\sup_{n:\,x\in V_n}u_{n,\eta}^{x}(y)\in\N\cup\{0,\infty\}.
\]
The auxiliary sink specifies the finite-volume approximation; there is no dissipative vertex at a finite site in the infinite lattice. Write
\[
 A_x(\eta):=\{y\in\Z^2:u_{\eta}^{x}(y)>0\},\qquad
 T_x(\eta):=\sum_{y\in\Z^2}u_{\eta}^{x}(y).
\]
Here an avalanche is called finite when \(T_x(\eta)<\infty\), equivalently when its toppled set is finite and every vertex topples finitely many times.

\textbf{Question.} For the intended infinite-volume stationary sandpile measure \(\nu\), is
\[
 \nu\bigl(\{\eta:T_x(\eta)<\infty\}\bigr)=1
\]
for each fixed addition site \(x\in\Z^2\)?

The supplied statement presents finiteness of the number of topplings at the addition site,
\[
 \nu\bigl(\{\eta:u_{\eta}^{x}(x)<\infty\}\bigr)=1,
\]
as a sufficient route to an affirmative answer. This sufficiency is retained as a source assertion, not established here. In particular, finiteness of a single toppling count does not by definition give finiteness of \(T_x\).

\textbf{Unresolved qualifications.} The source does not specify the construction of \(\nu\), the infinite-volume stabilization convention, or whether ``finite avalanche'' explicitly means finitely many topplings or finitely many toppled vertices. The definitions above give a conventional precise formulation for review; their agreement with the intended source formulation, and the asserted sufficient route, require confirmation.
\subsection{Short English statement}
After adding one grain at a fixed site of a stationary sandpile on \(\Z^2\), must the resulting avalanche contain only finitely many topplings with probability one? The source proposes proving that the addition site topples only finitely often as a sufficient route. The intended stationary measure and the justification of that route require review.
\subsection{Sources}
\begin{itemize}
\item Local supplied problem registry, permanent UnsolvedMath ID \(30003723\), title \emph{Finiteness of Two-Dimensional Sandpile Avalanches}, exact source problem code \texttt{catalogue entry}, in the \emph{Oberwolfach Reports Open Problems} collection. No individual proposer is supplied. The record gives proposed year 2018.
\item Supplied report citation: \emph{Strongly Correlated Random Interacting Processes} (2019), Oberwolfach Reports, pp.\ 187--253, DOI \texttt{10.4171/owr/2018/4}. Supplied source link: \url{https://doi.org/10.4171/owr/2018/4}. The record describes Oberwolfach Reports as documenting mathematical workshops and their open-problem discussions.
\item Supplied literature source, labelled ``original OWR report'': \url{https://doi.org/10.4171/OWR/2018/4}.
\item Supplied literature source, labelled ``Járai, sandpile models survey'': \url{https://arxiv.org/abs/1401.0354}.
\item Literature assessment dated exactly \texttt{2026-08-21}: the supplied curation labels the problem open and states that almost-sure avalanche finiteness after one grain addition on \(\Z^2\) remains open and that ``the cited implication'' has not been established. This is supplied dated curation, not an independently verified literature assessment. The background explicitly requests independent verification of time-sensitive status and cited claims.
\item The registry labels the statement \texttt{corrected\_verified} and reports manual semantic, typographic, and LaTeX review. These are supplied editorial metadata, not independent verification of the mathematical formulation or literature status.
\end{itemize}
\end{document}
